\section{Background}
\label{sec:background}

\paragraph{Flow Matching.}
Flow Matching (FM)~\cite{fm,rectified,stochasticflow} learns a velocity field that flows between a prior distribution and the data distribution. 
We consider the standard linear schedule $z_t=(1-t)\,x+t\,e$ with data $x\!\sim p_{\text{data}}$ and noise $e\!\sim p_{\text{prior}}$ (\eg Gaussian). 
Computing the time-derivative gives a \emph{conditional} velocity $v_c=e-x$.
Flow Matching learns a network $v_\theta$ to regress $v_c$ by minimizing a loss function in the $v$-space (namely, $v$-loss):
\begin{equation}
\label{eq:fm-loss}
\mathbb{E}_{t,x,e}\big\|\,v_\theta(z_t,t)-(e-x)\,\big\|^2.
\end{equation}
As one $z_t$ can be given by multiple pairs of $(x,e)$, the underlying unique regression target is the \textit{marginal velocity}~\cite{fm}:
\begin{equation}
\label{eq:marginal-v} 
v(z_t,t) \triangleq \mathbb{E}[\,v_c\mid z_t\,], 
\end{equation}
which is marginalized over all pairs $(x,e)$ that satisfy $z_t$ at $t$.
For brevity, we omit $t$ and denote $v(z_t,t)$ as $v(z_t)$, and $v_\theta(z_t,t)$ as $v_\theta(z_t)$, in the remaining of this paper.

At generation time, FM samples by solving an ODE: $\ddt{z_t}\,=\,v_\theta(z_t)$.
This is done by a numerical solver (\eg Euler or Heun) integrating from $t=1$ to $0$, with $z_1\sim p_{\text{prior}}$.

\paragraph{MeanFlow.}
By viewing $v(z_t)$ as the \textit{instantaneous} velocity field, MeanFlow (MF)~\cite{mf} introduces the \emph{average} velocity field between two time steps $r$ and $t$:
\begin{equation}
\label{eq:avg-vel}
u(z_t, r, t)\; \triangleq \;\frac{1}{t-r}\int_{r}^{t} v(z_\tau)\,d\tau.
\end{equation}
Again, for brevity, we omit $r$ and $t$ and simply denote $u(z_t, r, t)$ by $u(z_t)$. 
Directly integrating \cref{eq:avg-vel} at training time is intractable.
Instead, MF takes the derivative \wrt $t$ and obtains a \textit{MeanFlow identity}~\cite{mf}:
\begin{equation}
\label{eq:mf-identity}
u(z_t)\;=\;v(z_t)\;-\;(t-r)\,\frac{d}{dt}\,u(z_t),
\end{equation}
which is used to establish a feasible training objective.
The term $\ddt u$ is given by~\cite{mf}:
\begin{equation}
    \label{eq:total-deriv}
    \ddt u(z_t) = \partial_z u(z_t)\,v(z_t) + \partial_t u(z_t)
    \triangleq \mathtt{JVP}(u; v) .
\end{equation}
This can be computed by Jacobian-vector product (JVP), between the Jacobian $[\partial_z{u}, \partial_r{u}, \partial_t{u}]$ and a tangent vector $[v, 0, 1]$.
Here, for brevity, we introduce the notation $\mathtt{JVP}(u; v)$ for this JVP computed at $u(z_t)$ and $v(z_t)$.

MeanFlow \textit{parameterizes} average velocity by a network $u_\theta(z_t)$ (conditioned on $r$ and $t$, omitted in notation for brevity). This network is optimized to approximate the MeanFlow identity (\ref{eq:mf-identity}).
In the formulation of the original MF, \cref{eq:mf-identity} is implemented as:
\begin{equation}
\label{eq:u-tgt}
    \utgt = (e-x) - (t-r)\,\mathtt{JVP}(u_\theta; e-x),
\end{equation}
Here, two approximations are made~\cite{mf}: (i) the marginal $v(z_t)$ is replaced with the conditional $v_c=e-x$, same as Flow Matching; (ii) the true $u$ in $\mathtt{JVP}$ is replaced by its network prediction $u_\theta$.
With this target $\utgt$, MF optimizes: 
\begin{equation}
\label{eq:mf-loss}
\mathbb{E}_{t,r,x,e}\,\|u_\theta- \text{sg}(\utgt)\|^2,
\end{equation}
where ``sg'' denotes stop-gradient, which helps create an \emph{apparent} target for training.
Once trained, MF directly performs one-step sampling via $z_0 = z_1 - u_\theta(z_1)$ given $(r,t)=(0,1)$, with $z_1\sim p_{\text{prior}}$.